\documentclass[11pt]{article}
\usepackage[margin=1in]{geometry}
\usepackage[T1]{fontenc}
\usepackage[utf8]{inputenc}
\usepackage{lmodern}
\usepackage{booktabs}
\usepackage{graphicx}
\usepackage{tabularx}
\usepackage{amsmath}
\usepackage{microtype}
\usepackage[hidelinks]{hyperref}
\usepackage[numbers,sort&compress]{natbib}

\newcommand{\TODO}[1]{\textbf{[#1]}}

\title{Decision Models at the Game Table:\\
Where Non-Generative Classifiers Help and Fail\\
in LLM Tooling for Tabletop Role-Playing Campaigns}

\author{Kostadis Roussos\\
Independent researcher\\
\texttt{kostadis@gmail.com}}

\date{October 2026}

\begin{document}
\maketitle

\begin{abstract}
\emph{Decision models} are a recent class of non-generative language models: a backbone reads a
\emph{state} once and a small head returns a probability for every option of every typed question
(multiple choice, yes/no, ordinal score), with no text generated. We test six of them: TypeSafe's
hosted Jev, which introduced the SystemOne interface, and two open families released in October 2026, Cloudflare's
Clef and vLLM Semantic Router's Decision~2.0, served on a local NVIDIA DGX Spark (GB10). We ask
whether they can replace or assist a generative model in the judgment steps of a working LLM
pipeline for tabletop role-playing campaigns. Labels are the game master's own past rulings,
frozen into one versioned dataset: 766 reviewed transcript-correction cards, 864 triaged tokens,
497 spelling rulings and 115 entity-type merge decisions, plus 12 paired NPC-behaviour scenarios.
Decision models work where the label is a coarse property readable from meaning: Clef~27B and Jev
each forwarded about half of a transcript's unknown-token pile on a held-out split without
discarding any of 84 real names, and five of six models moved NPC behaviour in the expected
direction on 11 of 12 paired scenarios against 5 of 12 for a published dice-table system. They
fall short where the label is a string-level judgment (choosing a canonical spelling: 66--69\% for
the best decision models against 76\% for a local generative model), a house convention, or a
consistency property across many items (typing extracted facts one at a time created about as many
split entities as it repaired; a model-free per-entity vote matched the GM's filing 32 times in 33
where it decided). Jev is the exception that shows what to ask for: on spelling it abstained
instead of guessing, with 28 confident wrong names against 109 for Clef~27B. Fine-tuning a 4B decision model
on the GM's rulings recalibrated it (triage AUC 0.90 to 0.96) but did not teach the conventions. We
also report that an apparently strong verdict-prediction result was an artifact of the classifier
reading the upstream model's own hedging, and that vendor latency claims did not transfer to the
GB10 (127~ms measured vs.\ 18.4~ms claimed).
\end{abstract}

\section{Introduction}

Large language model pipelines for long-running creative work make many small judgment calls: is this
token a misheard proper noun, which entity does this fact describe, how would this character react.
In a pipeline that turns recorded game sessions into campaign records, each of these calls feeds
the next step, and an error compounds silently unless a human checkpoint catches it. A practitioner
therefore wants two things from the model making those calls: it should be right, and it should
\emph{know when it is not}, so that uncertain cases can be routed to the human instead of being
decided.

Decision models address both wants directly. Instead of generating an answer, they score a fixed set of options
and return a probability for each, in one forward pass. TypeSafe's \emph{Jev}~\cite{typesafe} is a
hosted model behind the \emph{SystemOne} request format. Two open families exposing the same format were released within days
of each other in October 2026: Cloudflare's \emph{Clef} (27B and 9B
backbones)~\cite{clef} and vLLM Semantic Router's \emph{Decision~2.0} (0.6B to 27B)~\cite{decision2},
both advertising low latency.

This paper reports a practitioner's evaluation of these models against real workloads rather than
benchmarks. The workloads come from an open-source campaign-management toolkit used to run several
tabletop role-playing campaigns; the labels come from one game master's (GM's) recorded rulings on
earlier model output. The open models ran on local hardware (two NVIDIA DGX Spark systems, GB10,
128~GB unified memory each), deliberately: part of the aim was to calibrate what local serving of
such models is like in practice. Jev, whose weights are not released, was called through its
hosted API as a reference point.

We ask four questions:
\begin{description}
  \item[RQ1] Can a decision model replace a generative model in the \emph{choice} steps of the
        pipeline (spelling selection, entity typing)?
  \item[RQ2] Can it act as a cheap \emph{filter} or \emph{flag} in front of an expensive model or a
        human reviewer?
  \item[RQ3] Are its probabilities useful for \emph{sampling} plausible behaviour (NPC reactions),
        where a distribution matters more than an argmax?
  \item[RQ4] Do the published speed claims hold on this hardware?
\end{description}

\paragraph{Contributions.} (i) Results on five tasks for six decision models and a generative
baseline, with labels taken from a GM's own rulings and frozen into a versioned dataset, including
two strong negative results; (ii) a characterization of \emph{which kinds of label} decision models
can and cannot learn from the input, and of the failure mode (abstain or guess) that separates
otherwise similar models; (iii) two methodological pitfalls that produced convincing but spurious
results in this study (label leakage through upstream prose; selection bias in a repair-only test
set); (iv) a fine-tuning test showing that the convention failure survives training on the GM's
own rulings; (v) deployment notes for the GB10, including a unified-memory accounting bug in an
upstream runtime.

\section{Background}

\subsection{Decision models and the SystemOne interface}
A request is a \texttt{state} (text) and a dictionary of named \texttt{questions}, each typed
\texttt{choice} (a dictionary of option keys to descriptions), \texttt{noul} (yes/no) or
\texttt{score} (an ordered list of levels). The model prefills the state and the question text once;
a small head scores every option of every question jointly. The response carries a probability per
option. TypeSafe defined the interface (\texttt{POST /v1/systemone}) for Jev, which it trains for
calibrated decisions rather than generated text and serves only as a hosted
API~\cite{typesafe}. Clef's release implements the head as a separate set of weights (``joint
schema head'') on Qwen-family backbones~\cite{clef}. Decision~2.0 is served by a standalone runtime
in vLLM Semantic Router~\cite{vllmsr,vllmsrvision} that exposes \texttt{POST /v1/decisions}, a
superset of SystemOne, with \texttt{/v1/systemone} as an alias~\cite{vllmsrruntime}. Because all
three share the interface, every experiment here ran against any model by changing only the
endpoint.

\subsection{The workloads}
The toolkit's pipeline turns a session's audio transcript into a reviewed transcript, a structured
summary, scene extractions and finally narrative prose. Three steps involve classification-shaped
judgments:
\begin{itemize}
  \item \textbf{Transcript spelling pass.} Automatic speech recognition mangles fantasy names
        (\emph{Zalthir} $\rightarrow$ ``Zalfir''). A pass collects unknown capitalised tokens,
        filters obvious non-names, proposes a canonical spelling for the rest, and asks the GM to
        approve each proposal on a review page (Approve / Reject / Discuss).
  \item \textbf{Ensemble fact extraction.} Five extraction passes over each chapter emit facts with a
        \texttt{type} (npc, monster, faction, location, object, event, thread, date). Downstream,
        facts are grouped into per-entity dossiers keyed by (type, canonical subject).
  \item \textbf{NPC behaviour at the table.} During play the GM must decide how a non-player character
        responds. A published system, FlexAI~\cite{flexai}, does this with d100 tables keyed on the
        creature's role and stance; it never sees the situation itself.
\end{itemize}

\subsection{The human-checkpoint principle}
The toolkit's design rule is that a model may draft, but \emph{precision decisions} (scope,
ordering, attribution, canonical spellings) require a human checkpoint before they feed another
model. This shapes what counts as a useful classifier: a model that cannot be trusted to decide may
still be valuable if it reliably \emph{orders} or \emph{flags} work for the human, provided nothing is
dropped silently.

\section{Setup}

\subsection{Models and serving}
Table~\ref{tab:models} lists the models. Clef and clef-flash were served by a thin FastAPI wrapper
around the release's own model code using PyTorch~2.11 (CUDA~13.0); the release's suggested vLLM
command would load only the chat backbone and silently drop the decision head. Decision~2.0 was served by
\texttt{vllm-sr-runtime}~\cite{vllmsrruntime} at a pinned commit, one model per container. Jev was
called as \texttt{jev-1.13.0}, pinned rather than the moving \texttt{jev-latest} alias, through
TypeSafe's public API from the same workstation that drove the local runs. The generative baseline
was \texttt{qwen3.8-flash-next} (NVFP4/FP8 hybrid quantisation, served by vLLM) with reasoning
either disabled (``fast'') or enabled with a 4{,}000-token budget (``thinking''). Questions,
option descriptions and prompts are identical across models.

\begin{table}[t]
\centering
\small
\begin{tabular}{lllr}
\toprule
Model & Family & Backbone & Median latency per request \\
\midrule
clef & Clef & Qwen3.8-27B & 420--960 ms \\
clef-flash & Clef & Qwen3.5-9B & 125--270 ms \\
Decision-2.0-Lux-9B & Decision 2.0 & Qwen3.5-9B & 127--295 ms (claimed 18.4 ms) \\
Decision-2.0-Nox-4B & Decision 2.0 & Qwen3.5 (4B) & 79--206 ms \\
Decision-2.0-Kai-0.6B & Decision 2.0 & Qwen3-0.6B & 17--53 ms (claimed 4.9 ms) \\
jev-1.13.0 & TypeSafe (hosted) & not disclosed & 116--121 ms$^\ast$ \\
qwen3.8-flash-next, fast & generative & --- & 2.0--4.9 s \\
qwen3.8-flash-next, thinking & generative & --- & \TODO{QT-LAT} s \\
\bottomrule
\end{tabular}
\caption{Models and measured median per-request latency for the request shapes used here (one
request, one to two questions, a few hundred to $\sim$2{,}000 input tokens). Ranges span tasks.
Local decision models were measured single-stream on the GB10; the generative model ran 8-way
(fast) or 6-way (thinking) concurrent. $^\ast$Jev's figure is end to end from the workstation,
including the internet round trip, and is not comparable to the GB10 rows. Claimed latencies are
from the model cards, measured on an unnamed GPU.}
\label{tab:models}
\end{table}

\subsection{Labels and the frozen dataset}
All labels are rulings the GM made on earlier model output in one campaign (an adapted
\emph{Out of the Abyss}). Four sources are used:
\begin{enumerate}
  \item \textbf{766 transcript-correction review cards} across 59 sessions: 662 approved, 98
        discussed, 4 ruled table chatter, 2 rejected.
  \item \textbf{864 triaged tokens} from the nine sessions whose review pages enumerated what was
        left alone: 628 left alone, 152 ruled table chatter, 84 that became approved
        name-correction cards. ``Left alone'' means the GM saw the token listed and did not pull it
        back, a weaker label than a card ruling.
  \item \textbf{497 spelling items}: 485 GM-approved name corrections whose canonical form exists
        in the campaign's name lists, plus 12 ruled table chatter (correct answer: leave as spoken).
  \item \textbf{115 entity-type merge rulings} over ensemble dossiers: 109 merged (with a GM-chosen
        primary type) and 9 kept separate (some groups mix outcomes).
\end{enumerate}
The original per-experiment scripts rebuilt their inputs from the live campaign tree on every run,
so the label sets grew while the study ran (the review cards from 525 to 766). For this paper every
input was rendered once, with each model-facing \texttt{state} already built, and frozen into a
read-only, versioned dataset; every model below was run against those same bytes, and a manifest
records each file's hash and the source commits. Rerunning on the frozen set reproduced the
deterministic results of the earlier, unfrozen runs exactly where the inputs had not changed (E4
per-fact typing, the E4 replay, E6 head-only triage). The campaign text is not released: it quotes
play sessions verbatim and the table-chatter labels name real people. The evaluation scripts and
the manifest are public~\cite{dgxfun}.

\subsection{Metrics}
For verdict prediction we report ROC AUC~\cite{hanley1982}. For filters we report, at each threshold,
the fraction of items forwarded and the number of real names \emph{missed}, since in this pipeline a
filtered-out real name is a silent deletion that nothing downstream can recover. For choice tasks we
report accuracy and separately count \emph{confident wrong names}, the error a tired reviewer
approves.

\section{Experiments and results}

\subsection{E1: Predicting the GM's verdict (and how it fooled us)}
\textbf{Task.} For each of the 766 review cards, ask ``Will the GM approve this proposed correction
as written?'' (\texttt{noul}). If this worked, the review page could be sorted so the GM looks
hardest at the likely rejections.

\textbf{Result.} On the cards as the upstream model wrote them, clef reached AUC~0.799 and Jev 0.807
(the other decision models 0.71--0.82). Rebuilding each card from raw sources only (transcript
lines, the second transcription, a canon lookup) raised clef to 0.919.

\textbf{Why it is not a result.} A one-feature baseline, ``does the card propose a fix at all'',
scores 0.897: 88 of the 123 cards with no proposed fix were not approved, against 16 of 643 cards
with one. Splitting by card type (Table~\ref{tab:e1}) shows that on open questions the strong score
existed only when the classifier could read the upstream model's prose, and fell \emph{below
chance} without it (clef 0.439, Jev 0.470). The classifiers were reading the drafting model's hedges (``could be X or
Y'', ``what was said?'') back to us. On proposals, the verdict is saturated (97.5\% approved) and the
16 non-approvals are too few to measure.

\begin{table}[t]
\centering
\small
\resizebox{\textwidth}{!}{\begin{tabular}{lrrrrr}
\toprule
& & \multicolumn{2}{c}{clef} & \multicolumn{2}{c}{jev} \\
\cmidrule(lr){3-4}\cmidrule(lr){5-6}
Card subset & $n$ (not approved) & upstream prose & raw sources & upstream prose & raw sources \\
\midrule
All cards & 766 (104) & 0.799 & 0.919 & 0.807 & 0.630 \\
Proposes a fix & 643 (16) & 0.549 & 0.722 & 0.418 & 0.689 \\
Open question, no fix & 123 (88) & 0.803 & 0.439 & 0.749 & 0.470 \\
\midrule
Baseline: ``has a fix'' & 766 (104) & \multicolumn{4}{c}{0.897} \\
\bottomrule
\end{tabular}}
\caption{E1, AUC for predicting the GM's verdict. The headline AUC comes from card type and from
upstream prose, not from judgment about the correction itself. The pattern holds for every decision
model tested (open-question AUC 0.73--0.81 with prose, 0.44--0.60 without).}
\label{tab:e1}
\end{table}

\subsection{E2: Choosing the canonical spelling}
\textbf{Task.} Given a transcript token, about 21 cues of raw transcript around it and the second
transcription's reading, pick the canonical spelling from a candidate list: the eight nearest names
by spelling and Double Metaphone~\cite{philips2000} similarity, plus \textsc{other} (a real name not
in the list) and \textsc{leave} (not a misheard name). Candidates came from canonical names only;
the glossary's wrong-form column was excluded because it contains these very rulings.

\textbf{Candidate recall caps every model.} The correct name was among the eight candidates for
385 of 485 items (79\%). The misses are mostly cases where neither spelling nor sound is close enough
to retrieve the name (\emph{Vira} $\rightarrow$ Ilvara, \emph{Peto} $\rightarrow$ Buppido), plus a
house convention of transcribing the GM's own name as ``GM''.

\textbf{Instruction sensitivity.} With the question phrased as ``which name was spoken'', both
models tried on that wording (clef-flash and the fast generative model) expanded a spoken surname into a full name (\emph{Misrim} $\rightarrow$ ``Ilvara Mizzrym''),
which would insert words nobody said. Rephrased as ``the correct spelling of exactly the words
spoken'', the generative model complied and clef did not.

\begin{table}[t]
\centering
\small
\resizebox{\textwidth}{!}{\begin{tabular}{lrrrrr}
\toprule
Model & Accuracy & Confident wrong names & Said \textsc{other} & Most-confident half & Median latency \\
\midrule
Decision-2.0-Kai-0.6B & 15.3\% & 54 & --- & 16.5\% & 0.05 s \\
Decision-2.0-Nox-4B & 19.7\% & 29 & --- & 20.2\% & 0.21 s \\
clef-flash & 30.0\% & 64 & --- & 32.7\% & 0.27 s \\
Decision-2.0-Lux-9B & 31.4\% & 44 & --- & 34.7\% & 0.30 s \\
clef & 66.4\% & 109 & 53 & 78.6\% & 0.96 s \\
jev & 69.0\% & 28 & 208 & 83.1\% & 0.12 s$^\ast$ \\
qwen, fast & 76.1\% & 87 & 72 & --- & 4.9 s \\
qwen, thinking & \TODO{QT-ACC} & \TODO{QT-WN} & --- & --- & \TODO{QT-LAT} s answered \\
\midrule
clef $\wedge$ jev agree & 90.3\% of 268 & 13 & --- & --- & --- \\
clef $\wedge$ qwen-thinking agree & \TODO{QT-AGREE} & \TODO{QT-AGREE-WN} & --- & --- & --- \\
\bottomrule
\end{tabular}}
\caption{E2, 497 items. ``Said \textsc{other}'' counts every \textsc{other} answer, right or wrong;
the Decision~2.0 models and clef-flash mostly answered \textsc{leave} or \textsc{other} instead of
reading the candidates. ``Most-confident half'' is the accuracy of the half of items with the
highest reported confidence; the fast generative model's stated confidence did not separate them.
\TODO{QT-NOTE}}
\label{tab:e2}
\end{table}

\textbf{Result (Table~\ref{tab:e2}).} Decision models lose to the generative baseline, and the open
models below 27B are not usable for this task. The two strongest decision models reach similar
accuracy by different routes. Clef commits: 109 of its 167 errors are a wrong name, the error a
reviewer approves. Jev abstains: 119 of its 154 errors are \textsc{other} on an item whose name was
in the list, and it answered \textsc{other} correctly for 89 of the 100 items whose name was not.
Its confidence is informative: the 79 items it answered with confidence $\geq 0.9$ were 92.4\%
correct with 2 wrong names, a slice that could be pre-filled for the GM. Note the test set's bias:
every item is a proposal the upstream model made and the GM approved, so it measures how close these
models come to approved work and cannot reveal the upstream model's own misses.

\subsection{E3: Triage of the unknown-token pile}
\textbf{Task.} Before any expensive judgment, classify each capitalised unknown token as filler,
rules term, real-world chatter, or campaign name (\texttt{choice}), and forward it only if
$P(\text{campaign})$ clears a threshold. Nothing is deleted: tokens below the threshold still appear
in the review page's ``left alone'' list.

\begin{table}[t]
\centering
\small
\resizebox{\textwidth}{!}{\begin{tabular}{lrrrrrr}
\toprule
Model & AUC & Fwd., 0 of 84 missed & Fwd., $\leq$2 missed & Held-out halves (misses) & Chatter recognised & Median \\
\midrule
clef & 0.933 & 34--42\% & 27--39\% & 43\% (0), 49\% (0) & 136/152 & 419 ms \\
jev & 0.931 & 42--55\% & 37--50\% & 51\% (0), 47\% (0) & 143/152 & 116 ms$^\ast$ \\
clef-flash & 0.904 & 92--95\% & 47--58\% & 90\% (1), 100\% (0) & 128/152 & 124 ms \\
Decision-2.0-Lux-9B & 0.896 & 75--89\% & 33--49\% & 69\% (1), 96\% (0) & 129/152 & 127 ms \\
Decision-2.0-Nox-4B & 0.900 & 82--89\% & 34--46\% & 96\% (0), 94\% (0) & 118/152 & 79 ms \\
Decision-2.0-Kai-0.6B & 0.796 & 97--98\% & 49--55\% & 97\% (0), 100\% (0) & 7/152 & 17 ms \\
qwen, fast & --- & --- & 24\% (6 missed) & --- & 145/152 & 2.0 s \\
\bottomrule
\end{tabular}}
\caption{E3, 864 tokens, 84 of them real names. ``Fwd.'' is the share of tokens sent on to the
expensive step at the threshold achieving the stated number of misses; ranges span two readouts of
the same answer ($P(\text{campaign})$ and $1 - P(\text{filler}) - P(\text{realworld})$). ``Held-out
halves'': a threshold set at half the lowest name score on one half of the sessions, applied to the
other half.}
\label{tab:e3}
\end{table}

\textbf{Result (Table~\ref{tab:e3}).} This is the task where decision models earned a place. Clef~27B
and Jev each roughly halve the expensive step's input without a single silent loss, and both survive a
held-out split; Jev also recognised the most table chatter of any decision model. Smaller open models trade that guarantee for speed; the names they drop recur across
models (e.g.\ \emph{Heel Strike}).

\subsection{E4: Entity typing in ensemble extraction}
\textbf{Per entity.} For each of the 92 usable merged groups, the classifier saw a sample of the
entity's raw facts with the original type labels hidden. A free majority vote over the extraction
passes' own types matched the GM's primary on 82/92; the decision models matched 79--82 (clef and
Nox 82, Jev 80) and the fast generative model 80. Voting with the model deciding only when no type
holds 60\% of facts scored 85--87 for every model (Jev 86). Most disagreements are filing conventions that the facts do not reveal:
the GM filed the demon lords Zuggtmoy, Juiblex and Demogorgon as \emph{monster} but Yeenoghu, also a
demon lord, as \emph{npc}; the campaign's own entity registry types all four \emph{npc}. These are canon questions for the GM, not model errors.

\textbf{Per fact.} The idea tested was to let the generative model extract facts and a decision model
assign each fact's type. On the 1{,}691 facts of those entities, extraction had typed 266 with a
different entity type, splitting 78 entities. Re-typing reduced splits (Table~\ref{tab:e4}), but this
set was \emph{selected on extraction's mistakes}, so any second opinion looks good on it. On a control
set of 54 entities that extraction had typed consistently (839 facts), every re-typer created new
splits, and clef moved 21\% of facts into non-entity types (\emph{event}, \emph{thread}, \emph{date}),
removing them from their dossiers.

\begin{table}[t]
\centering
\small
\resizebox{\textwidth}{!}{\begin{tabular}{lrrrr}
\toprule
Who types each fact & Per-fact accuracy & Split entities (of 92 merged) & New splits (of 54 consistent) & Facts moved to non-entity types \\
\midrule
extraction (as is) & 84.3\% & 78 & 0 (by selection) & --- \\
clef & 76.8\% & 29 & 7 & 255 + 177 \\
jev & 84.2\% & 35 & 9 & 87 + 80 \\
clef-flash & 81.6\% & 46 & 13 & 120 + 68 \\
Decision-2.0-Lux-9B & 86.3\% & 41 & 12 & 22 + 22 \\
Decision-2.0-Nox-4B & 83.6\% & 44 & 12 & 42 + 28 \\
Decision-2.0-Kai-0.6B & 63.6\% & 65 & 34 & 200 + 137 \\
qwen, fast & 88.3\% & 38 & 12 & 28 + 17 \\
\bottomrule
\end{tabular}}
\caption{E4, per-fact re-typing. Accuracy is on the repair set against the GM's primary type.
``Facts moved'' gives repair set + control set. Projected over the campaign's $\sim$470 consistent
dossiers, re-typing creates roughly as many splits as it repairs (clef, Jev) or more (the others).}
\label{tab:e4}
\end{table}

\textbf{As an extra signal.} Adding decision-model votes to the majority vote hurt (81/92 voting over
clef's per-fact types, 78/92 over Jev's; 82--84/92 pooled with extraction's; 84/92 for extraction
alone on this fact set). Used instead as a \emph{second opinion} that flags subjects where it
disagrees with the vote, clef flagged 15/92 and caught 8 of the vote's 10 errors (7 false alarms);
Jev (15 flags, 7/10), the fast generative model (17, 8/10) and clef-flash (18, 8/10) did about as
well.

\textbf{Interpretation.} The failure is structural. Per-fact typing gives an entity with dozens of
facts dozens of independent chances to draw one odd label, and one odd label is a split. A
deterministic model guarantees the same answer for the same input, but every fact is a different
input. The repair is to type once per entity, which needs no model.

\textbf{Replaying entity-level typing.} We then specified that repair as a resolution chain run
before dossiers are built: (1) GM filing conventions, (2) the GM's prior merge rulings, (3) the entity
registry's type, (4) a majority vote over the subject's own facts if one type holds at least 60\%,
otherwise (5) a GM review queue. No step calls a model. We replayed it, read-only, against the real
bundling rules on the 438 subjects of the campaign's extraction corpus that have at least one
\emph{monster}-typed fact (50 of them named), scoring against the 41 GM merge rulings whose files fall
in that corpus (Table~\ref{tab:e4r}). The informative arm is the vote alone, because it uses no GM
data: it is what a new, unreviewed chapter gets. When it decides, it matches the GM's primary type 32
times in 33; the one miss is a subject the GM had already flagged as mis-attributed. What it declines
to decide is the right set to show the GM: the nine queued named subjects are mostly the demon-lord,
deity and place-versus-organisation questions above. A threshold sweep put the knee at 60\%: 50\% adds
three wrong types, 67\% more than doubles the queue. Adding the GM's merge rulings reproduces all 40
merged groups that fell in the slice (partly circular, since it is scored against those rulings).

\begin{table}[t]
\centering
\small
\resizebox{\textwidth}{!}{\begin{tabular}{lrrr}
\toprule
 & Per-fact types (today) & Vote only, $\geq$60\% & GM rulings + vote \\
\midrule
Named subjects split across types (of 50) & 43 & 9 & 1$^\dagger$ \\
GM groups in one dossier under the GM's primary (of 41) & 0 & 32 & 40 \\
GM groups in one dossier under the wrong type & 1 & 1 & 0 \\
Named subjects sent to the GM queue & --- & 9 & 0 \\
Dossiers & 667 & 592 & 577 \\
\bottomrule
\end{tabular}}
\caption{E4, entity-level typing replayed on the \emph{monster} slice. $^\dagger$The remaining split is
one the GM ruled: a group whose \emph{npc} and \emph{monster} files were merged while its
\emph{faction} file was kept separate.}
\label{tab:e4r}
\end{table}

The replay also surfaced two defects in the specification and one in the pipeline. Keep-separate
rulings must apply to a \emph{file} (type and subject), not a subject, because a single group can be
partly merged and partly kept apart. The review queue must hold named subjects only; all 20 subjects
queued once the GM's rulings were applied were generic creatures. And subjects are canonicalised only
through the registry's aliases, so misspellings known to the spelling glossary but not the registry,
and extractions made before a source-text correction, split subjects before typing ever runs. The
underlying flaw is general: the pipeline asks the generative model for precision decisions (a fact's
type decides its scope, its subject decides its identity) and then stores the GM's corrections where
the pipeline never reads them, so every run rebuilds what the GM already fixed.

\subsection{E5: NPC behaviour}
\textbf{Task.} Twelve paired combat scenarios, each pair differing in one fact, with the expected
direction stated before running (e.g.\ a goblin with all allies dead and an open tunnel behind it
should flee more; the same goblin in a dead-end alcove should flee less; a mindless zombie should not
flee). Two questions per request in one call: what the creature does (FlexAI's six outcomes) and whom
it targets (FlexAI's eight targeting rules). FlexAI's own d100 tables are the baseline; they see only
role and stance.

\begin{table}[t]
\centering
\small
\resizebox{\textwidth}{!}{\begin{tabular}{lrr}
\toprule
Policy & Pairs moving the expected way ($\Delta > 0.02$) & Goblin $P(\text{flee})$: unhurt $\rightarrow$ alone $\rightarrow$ cornered \\
\midrule
FlexAI d100 tables (role $\times$ stance) & 5 / 12 & 0.15 $\rightarrow$ 0.16 $\rightarrow$ 0.00 \\
clef & 11 / 12 & 0.03 $\rightarrow$ 0.75 $\rightarrow$ 0.43 \\
clef-flash & 11 / 12 & 0.03 $\rightarrow$ 0.48 $\rightarrow$ 0.09 \\
Decision-2.0-Lux-9B & 11 / 12 & 0.01 $\rightarrow$ 0.59 $\rightarrow$ 0.17 \\
Decision-2.0-Nox-4B & 11 / 12 & 0.01 $\rightarrow$ 0.49 $\rightarrow$ 0.10 \\
Decision-2.0-Kai-0.6B & 5 / 12 & 0.05 $\rightarrow$ 0.16 $\rightarrow$ 0.20 \\
jev & 11 / 12 & 0.00 $\rightarrow$ 0.89 $\rightarrow$ 0.09 \\
qwen, fast (asked for a distribution) & 9 / 12 & 0.00 $\rightarrow$ 0.70 $\rightarrow$ 0.00 \\
\bottomrule
\end{tabular}}
\caption{E5. Two pairs are imperfect: in one, clef's answer (an archer engaged in melee draws its
sword, \emph{attack secondary} 0.64) is arguably better than the expected \emph{maneuver}; in
another, the shared party description contradicted the scenario's premise, invalidating the pair for
all policies. The scenarios and expectations were written by the experimenter.}
\label{tab:e5}
\end{table}

\textbf{Result (Table~\ref{tab:e5}).} Situation-aware decision models clearly beat the dice tables,
and their probabilities are graded, which is what a sampler needs. They differ in how far they move:
clef keeps a cornered goblin at 0.43 \emph{flee}, Jev drops it to 0.09 after putting 0.89 on fleeing
when the goblin is alone with an open tunnel. Which is right is a table-style question, not one the pairs can
settle. The generative model's stated
distributions were nearly degenerate (many 0.00 and 1.00 values, e.g.\ 1.00 on a single target), so
sampling from them adds little variety. On this evidence we built a table-side prototype: the GM types
one line describing the situation, a decision model returns stance probabilities over eight social
responses and samples one, and the generative model voices a line in that stance. With
Decision-2.0-Nox-4B a suggestion takes 0.25--0.8~s including the NPC's dossier, and the line about
2.5--3~s. A labelled evaluation of social responses against 60 past moments has been prepared but
not yet adjudicated by the GM; we do not report it.

\subsection{E6: Fine-tuning a decision model on the GM's rulings}
\textbf{Question.} Can the failures above be trained away? We fine-tuned Decision-2.0-Nox-4B on
the GM's rulings for two tasks: triage (E3) and per-fact entity typing (E4, where the labels encode
the GM's filing conventions). Jev cannot be tried here: TypeSafe does not fine-tune it per
customer and serves the same weights to every account~\cite{typesafe}.

\textbf{Method.} The released Transformers wrapper refuses \texttt{train()}, but the runtime's
native backbone and head are ordinary PyTorch modules. We loaded the model through the runtime
exactly as the server does and confirmed that cached features reproduced the served probabilities
(mean $|\Delta p| = 0.0009$). We then trained in two stages: \emph{head only} on cached features,
and a \emph{rank-16 LoRA} on the MLP and attention projections plus the head, with gradients through
the backbone. The loss is valid-$k$ cross-entropy ($-\log\sum_{k\in\text{valid}}p_k$, the loss
family the release reports), with an L2 pull toward the released head. Evaluation is
cross-validated by \emph{session} for triage (3 folds) and by \emph{entity} for typing (5 folds),
so no test entity or session is seen in training.

\begin{table}[t]
\centering
\small
\resizebox{\textwidth}{!}{\begin{tabular}{lrrrrrr}
\toprule
& \multicolumn{3}{c}{Triage (864 tokens)} & \multicolumn{3}{c}{Entity typing (2{,}530 facts, 146 entities)} \\
\cmidrule(lr){2-4}\cmidrule(lr){5-7}
Nox-4B & AUC & Fwd.\ @ 0 miss & Fwd.\ @ $\leq$2 & Per-fact acc. & Split (92 / 54) & Entity majority \\
\midrule
released & 0.900 & 82\% & 34\% & 87.2\% & 44 / 12 & \TODO{ENT-REL} / 146 \\
head only & \textbf{0.963} & \textbf{67\%} & 34\% & \textbf{91.3\%} & 34 / 9 & --- \\
LoRA r16 + head & 0.935 & 83\% & 36\% & \TODO{LT-ACC} & \TODO{LT-SPLIT} & \TODO{ENT-LORA} / 146 \\
\bottomrule
\end{tabular}}
\caption{E6, held-out results. ``Fwd.'' is the share of tokens forwarded at the best held-out
threshold for the stated number of missed names. ``Split'' counts entities whose facts land in more
than one entity type (repair set / control set). Head-only settings were chosen from four
candidates on the same cross-validation, so its rows are mildly optimistic. LoRA training is not
bit-reproducible; an earlier run on the unfrozen inputs gave triage AUC 0.915.}
\label{tab:e7}
\end{table}

\textbf{Result (Table~\ref{tab:e7}).} Head-only tuning \emph{recalibrates} the readout. Triage
AUC rose from 0.900 to 0.963, and per-fact typing accuracy from 87.2\% to 91.3\%, with almost no
facts drifting into non-entity types (42 + 28 to 0 + 1). It did not learn the conventions. The demon lords the GM files
as \emph{monster} moved \emph{toward} \emph{npc} (Demogorgon: 12 of 12 facts \emph{monster} released, 1 of 12
tuned; Juiblex: 7 of 14 to 0), outvoted by a label distribution that is 64\% \emph{npc} and by the GM's own
contrary ruling on Yeenoghu. In the earlier, unfrozen run, weak regularisation overfit badly: a threshold chosen on the training
sessions missed 19 held-out names. \TODO{LORA-RESULT} Three triage names
(\emph{Heel Strike}, \emph{Zoom}, \emph{Pick / Shine}) were missed by every model and every
tuning, which suggests a limit of the label given this input rather than of model capacity.

\subsection{E7: Serving on the GB10}
\textbf{Latency.} Table~\ref{tab:models} summarises it. Decision~2.0's runtime ships custom kernels only for
AMD's gfx942; on CUDA it uses pure-PyTorch reference kernels, and its CUDA path is marked unvalidated
upstream. Measured medians on the GB10 were 3--7$\times$ the model-card figures (Kai 17 vs.\ 4.9~ms; Lux 127 vs.\ 18.4~ms). For Clef, the GB10 was
3.5$\times$ slower than Cloudflare's published H200 medians for the 9B model but only 1.9$\times$ for the
27B, consistent with fixed per-request overhead dominating small models. Hosted Jev answered in
119~ms at the median and 232~ms at the 99th percentile, internet round trip included; a few
requests stalled for tens of seconds in rate-limit back-off. The whole Jev run (6{,}292 requests,
4.2M input tokens) cost about \$0.18 at the listed price.

\textbf{Correctness.} With no upstream golden check on CUDA, we compared Kai-0.6B's CUDA and CPU
outputs on identical requests: probabilities agreed within 0.005.

\textbf{Memory.} The runtime's placement check used \texttt{torch.cuda.mem\_get\_info}, which on the
GB10's unified memory excludes reclaimable page cache; the box reported 9--18~GiB ``free'' with 34~GB
available, and reading the weights to verify checksums refilled the cache just before the check. We
patched the check to use the kernel's \texttt{MemAvailable} on integrated devices. Separately, Lux-9B
occupied about 50~GB (33~GB resident in the process) for 16~GB of BF16 weights; co-hosted with Clef the
system fell to 1~GB available and the runtime crashed once. For the frozen-dataset rerun each family
was therefore served alone, one after another.

\section{Discussion}

\paragraph{What kind of label is it?} Across the five tasks a single distinction predicted success.
Decision models succeeded when the label is a \emph{coarse property readable from meaning}: is this
token chatter or a campaign name; would a cornered creature flee. They fell short when the label is
(a) a \emph{string-level} judgment (which spelling, exactly which words), (b) a \emph{convention}
the text does not reveal (how this GM files demon lords), or (c) a \emph{consistency property across
items} (one type per entity across many facts). Most judgment steps in this pipeline are of the second
kind, which is why the decision models' role ended up narrow. Fine-tuning (E6) did not move this
boundary: training on a few thousand of the GM's own rulings sharpened the readout on clear cases,
but a sparse, partly contradictory convention was outvoted rather than learned. Jev, a closed
model trained by its vendor for calibrated decisions, sits on the same side of every
boundary as the open models; it differs in how it fails.

\paragraph{Abstaining is the property to ask for.} On spelling, clef and Jev were equally accurate,
but clef's errors were mostly wrong names and Jev's were mostly abstentions, and Jev's
confidence separated right from wrong more sharply (its answers at confidence $\geq 0.9$: 79 items,
92.4\% correct; clef's: 20 items, 70\%). For a pipeline with a human
checkpoint that difference is decisive: an abstention costs the GM a lookup, a confident wrong name
costs a silent error in the record. Accuracy tables hide it; we recommend counting the two kinds of
error separately for any decision model placed in front of a reviewer.

\paragraph{Filters and flags, not deciders.} Every useful configuration kept the human checkpoint
intact: triage that orders work without deleting it, a second opinion that adds subjects to a review
queue, and sampled NPC suggestions the GM reads. None of these removes a decision from the human.

\paragraph{A ruling should be made once.} The checkpoint only pays if its output flows back in. The
ensemble pipeline asked a generative model for per-fact precision decisions and stored the GM's 109
corrective merges where no stage read them, so each run recreated the same splits for the GM to fix
again. Resolving each precision field against the GM's accumulated rulings, with a deterministic
default and a queue for the rest, did more for consistency than any classifier we tried (E4 replay).
The better model was not the missing piece; the missing piece was reading the rulings back.

\paragraph{Calibration has two uses.} Graded probabilities mattered in two different ways: to set a
zero-loss filter threshold (E3) and to sample plausible behaviour (E5). The generative baseline's
self-reported probabilities served neither purpose.

\paragraph{Local versus hosted.} On the two tasks where decision models earned a place, the hosted
Jev matched the best local model (clef~27B) on triage and every 4B-and-larger local model on NPC
behaviour, at a fraction of the latency of clef~27B. The local path's value here was not quality: it
was the ability to fine-tune (E6), to inspect the head, and to find the GB10-specific failures (E7).

\paragraph{Two ways a convincing result was wrong.} E1's AUC came from the classifier reading the
upstream model's prose, and E4's per-fact improvement came from a test set selected on errors. Both
looked strong until a baseline or a control was added. We recommend both checks for any evaluation
built from a pipeline's own review records, and freezing the inputs before comparing models: our
own label sets grew by half between the first runs and this paper.

\section{Limitations}
This is a single-practitioner study: one campaign, one GM, one hardware platform. Several label sets
are small (84 real names in E3; 16 non-approved proposals in E1; 12 scenario pairs in E5). Some labels
are weak (``left alone'' in E3), and some test material was written by the experimenter (E5). The E2
set was built from approved proposals and cannot reveal the upstream model's misses. Prompts and
option descriptions were written once and not tuned, which is fair to a drop-in classifier but
unfavourable to it; TypeSafe's own guidance favours structured instructions and criteria, which we
did not use, so Jev in particular may be under-served. Jev is a hosted model of undisclosed size
and training data, called at one pinned version; its results describe that version, and its
latency includes our internet path. Running it sent the frozen inputs, which include transcript
excerpts, to a third-party service. We did not verify that the locally served Clef matches Cloudflare's hosted
service. Decision~2.0 ran on a runtime path its authors mark as unvalidated on CUDA; our own CPU/GPU
parity check covers only the smallest model. Local latencies are single-stream; the runtime's batching
profiles were not tested. Fine-tuning explored one LoRA configuration (rank 16, one or two epochs) and
selected head-only settings on the same cross-validation folds. The entity-level typing replay
covers only subjects with a \emph{monster}-typed fact in one of two extraction corpora, and 41 of the
64 relevant GM rulings; the other entity types are not yet replayed.

\section{Conclusion}
Decision models are a useful addition to a local LLM toolkit, in a narrow role: fast, calibrated
classification of coarse properties, placed where it orders or flags work rather than deciding it.
In this pipeline that meant triaging transcript tokens without silent losses and generating
situation-aware NPC behaviour at the table. They did not replace a generative model for spelling
decisions, could not infer house conventions, and cannot impose consistency across items
that are judged one at a time; that consistency came instead from typing once per entity and feeding
the GM's own rulings back into the pipeline. Where they fall short, the useful question is whether a
model abstains or guesses: the hosted Jev mostly abstained and the open models mostly guessed. On
the GB10, the smaller Decision~2.0 models give most of the behavioural quality at a fraction of the
memory, but published latency claims should be re-measured on the target hardware.

\section*{Data and code availability}
Serving configurations, the GB10 patch, all evaluation scripts, and the frozen dataset's manifest
(file hashes, counts and source commits) are available in the public repository~\cite{dgxfun}. The
campaign data itself is not released because it quotes play sessions verbatim and names real people.

\section*{Use of AI assistance}
The experiments, analysis and a first draft of this paper were carried out by an AI coding agent
(Claude, Anthropic) working under the author's direction, with the author reviewing results and
making the decisions recorded in the paper. The author takes full responsibility for the content.

\bibliographystyle{plainnat}
\begin{thebibliography}{10}

\bibitem[TypeSafe(2026)]{typesafe}
TypeSafe.
\newblock Jev and the System One API.
\newblock Documentation, \url{https://docs.typesafe.ai}, model \texttt{jev-1.13.0}, accessed October 2026.

\bibitem[Cloudflare(2026)]{clef}
Cloudflare.
\newblock Clef and Clef-flash decision models.
\newblock Model cards, \url{https://huggingface.co/Cloudflare/clef} and
  \url{https://huggingface.co/Cloudflare/clef-flash}, October 2026.

\bibitem[vLLM Semantic Router Team(2026a)]{decision2}
vLLM Semantic Router Team.
\newblock Decision 2.0: towards open foundation decision models.
\newblock Model collection, \url{https://huggingface.co/collections/vllm-sr/decision-20},
  October 2026.

\bibitem[vLLM Semantic Router Team(2026b)]{vllmsrruntime}
vLLM Semantic Router Team.
\newblock Built-in model runtime, phase 1: serve Decision 2.0 from vllm-sr.
\newblock GitHub issue \#4479, \url{https://github.com/vllm-project/semantic-router/issues/4479}, 2026.

\bibitem[vLLM Semantic Router(2025)]{vllmsr}
vLLM Semantic Router.
\newblock When to reason: Semantic router for vLLM.
\newblock arXiv:2510.08731, 2025.

\bibitem[vLLM Semantic Router(2026c)]{vllmsrvision}
vLLM Semantic Router.
\newblock The workload-router-pool architecture for LLM inference optimization: A vision paper from
  the vLLM Semantic Router project.
\newblock arXiv:2603.21354, 2026.

\bibitem[Infinium Game Studios(2020)]{flexai}
Infinium Game Studios.
\newblock \emph{FlexAI Guidebook}.
\newblock 2020.

\bibitem[Hanley and McNeil(1982)]{hanley1982}
J.~A. Hanley and B.~J. McNeil.
\newblock The meaning and use of the area under a receiver operating characteristic (ROC) curve.
\newblock \emph{Radiology}, 143(1):29--36, 1982.

\bibitem[Philips(2000)]{philips2000}
L.~Philips.
\newblock The double metaphone search algorithm.
\newblock \emph{C/C++ Users Journal}, 18(6):38--43, 2000.

\bibitem[Roussos(2026)]{dgxfun}
K.~Roussos.
\newblock dgx-fun: experiments with local LLM serving on NVIDIA DGX Spark.
\newblock \url{https://github.com/kostadis/dgx-fun}, 2026.

\end{thebibliography}

\end{document}
